\chapter{Web APIs}

\section{Making Web Requests with HTTP}

When you open a weather app, it does not know the forecast. It sends a
short message across the internet to a server, asking ``What is the
weather in Denver?'', and the server sends back an answer. Your phone does
this hundreds of times a day. The How the Internet Works chapter showed how
those messages travel as packets; in this chapter, you will learn the language
the messages are written in, and then build a server of your own.

The Hypertext Transfer Protocol (HTTP) is the protocol used for
transmitting hypertext over the World Wide Web. It is the foundation
of any data exchange on the web, and it is a protocol used for
transmitting hypertext requests from clients (like a user's browser)
to servers, which respond with the requested resources.\index{HTTP}

\subsection{HTTP Requests}\index{HTTP request}

An HTTP request is made up of several components:

\begin{itemize}
    \item \textbf{Method}: The HTTP method, like GET (retrieve data), POST (send data), PUT (update data), DELETE (remove data), and so on.
    \item \textbf{URL}: The URL of the resource to retrieve, send data to, update or delete.
    \item \textbf{Headers}: Additional information about the request or response, like the content type of the body.
    \item \textbf{Body}: The body of the request, used when sending data in POST or PUT requests.
\end{itemize}

One way to think about an HTTP request is to imagine ordering food at
a restaurant. The method is the type of action you want to take, such
as asking to see the menu, placing an order, changing an order, or
canceling it. The URL is like telling the server exactly which
restaurant counter to go to. The headers are like the extra details
you give, such as saying you want vegetarian food or that you would
like water with your meal. The body is the actual order information,
such as the list of food you want to buy.

The method tells the server what kind of action the client wants to
perform. For example, a GET request is used when the client wants to
receive information from the server, such as loading a web page or
asking for a list of items. A POST request is used when the client
wants to send new information to the server, such as creating an
account or submitting a form. A PUT request is used to update
something that already exists, and a DELETE request is used to remove
something.

The URL works like an address. It tells the request where to go. If
the method explains what the client wants to do, the URL explains
where that action should happen. For example, if a web app wants a
list of users, it might send a request to a URL that ends in
\verb|/users|. If it wants information for one specific user, it might
send a request to a URL like \verb|/users/5|.

Headers help the client and server understand each other better. They
carry extra information along with the request or response, such as
what type of data is being sent. This is similar to writing
``Fragile'' on a package or adding a note that says ``Please handle
with care.'' The message is still the same package, but the note helps
explain how it should be handled.

The body is the part of the request that contains the actual data
being sent. For example, when a user fills out a sign-up form, the
name, email address, and password they enter may be placed in the body
of the request. Not every request needs a body. GET requests usually
do not send one, because they are mainly asking to receive information
rather than send it.


\subsection{Using HTTP with Web-Based APIs}


Software developers often use HTTP to interact with web-based APIs. An
Application Programming Interface (API) is a set of rules that allows
programs to talk to each other. The developer creates the API on the
server and allows the client to talk to it.

When a developer makes a request to an API endpoint, they are asking
the server to either send them some data or receive some data from
them. The response from the server will often be in a format like JSON
or XML, which the developer can then use in their own application.

For example, a developer might make a GET request to
\texttt{https://api.example.com/users} to retrieve a list of all
users. The server would respond with a list of users in a format like
JSON.

An \index{API endpoint}API endpoint is a specific URL that a client can use to communicate
with a server. Each endpoint usually has a clear purpose. One endpoint
might return a list of books, while another might return information
about a single book. In the same way that a restaurant has different
counters for ordering food, paying, or picking up an order, an API can
have different endpoints for different tasks.

When a client sends a request to an API, the server reads the request
and decides how to respond. If the request is valid, the server sends
back a response with the requested data or a message saying the action
worked. If something is wrong, the server sends back a response
explaining that there was a problem. This back-and-forth communication
is one of the main ways web applications work behind the scenes.

For example, imagine a school website that stores a list of students.
A client could send a GET request to an endpoint like \texttt{/students}
to ask for the full list. It could send a POST request to the same
endpoint to add a new student. It could send a PUT request to
\texttt{/students/3} to update the information for one student, or a
DELETE request to \texttt{/students/3} to remove that student from the
list. These requests follow the same HTTP ideas, even if the website
or app itself looks very different on the screen.

Many modern apps work this way. A weather app may send a request to
get the latest forecast. A music app may send a request to get a list
of songs. A shopping app may send a request to place an order. Even
though users may only see buttons, menus, and pictures, HTTP requests
are often moving data in the background.


\subsection{HTTP Responses}\index{HTTP response}

After a server receives an HTTP request, it sends back an HTTP
response. The response is the server's answer to the client. Just as a
request has several parts, a response also has its own important
parts.

An HTTP response usually includes a status code, headers, and
sometimes a body. The status code tells the client what happened. The
headers give extra information about the response. The body contains
the returned data, if there is any.

A simple way to think about a response is to imagine asking a librarian
for a book. The librarian might say, ``Here is the book you asked
for,'' ``That book is checked out,'' or ``We do not have that
book.'' These are different kinds of responses to the same kind of
request.

Some common \index{status code}HTTP status codes are:

\begin{itemize}
    \item 200: The request worked, and the server is sending back the requested data.
    \item 201: The request worked, and the server created something new.
    \item 204: The request worked, and there is nothing to send back (common after a delete).
    \item 401: The client has not proven who it is (for example, a missing API key).
    \item 400: The request had a problem, such as missing or incorrect information.
    \item 404: The server could not find what was requested.
    \item 500: Something went wrong on the server.
\end{itemize}

These codes help the client understand the result quickly. Instead of
reading a long message first, the client can look at the status code
and get a fast summary of what happened.


\subsection{Data Formats in Responses}

When a server sends data back to a client, that data needs to be
organized in a way both sides can understand. Two common formats are
JSON and XML. These formats are simply ways of arranging information
so that programs can read it clearly and consistently.

\index{JSON}JSON is one of the most common formats used with web APIs. It is
popular because it is lightweight and easy for both humans and
programs to read. A JSON response might include data written as
key--value pairs. For example, a user might be described with data
such as a name, age, and email address.

Here is a small example of JSON:

\begin{verbatim}
{
  "name": "Ava",
  "grade": 10,
  "club": "Robotics"
}
\end{verbatim}

In this example, \texttt{"name"}, \texttt{"grade"}, and \texttt{"club"}
are the labels, and \texttt{"Ava"}, \texttt{10}, and
\texttt{"Robotics"} are the values. This is similar to filling out an
information card. Each label tells what kind of information is being
stored, and each value is the actual information.

XML is another format that may be used. It organizes information with
opening and closing tags. While it is less common in many modern
beginner examples, it is still useful to know that different APIs may
return data in different formats.


\section{Calling a Real API from Python}\index{API!calling from Python}

Enough theory: let us ask a real server for real data. Open-Meteo is a free weather service with a public API
that needs no key. You will send it a GET request from Python and get today's weather back as JSON.

Python's \verb|requests| library makes HTTP requests easy. Install it the same way you installed matplotlib:
\begin{minted}{shell}
pip install requests
\end{minted}

Then create a file called \filename{weather.py}:

\begin{minted}{python}
import requests

url = "https://api.open-meteo.com/v1/forecast"
params = {
    "latitude": 39.74,       # Denver
    "longitude": -104.99,
    "current": "temperature_2m,wind_speed_10m",
    "temperature_unit": "fahrenheit",
    "wind_speed_unit": "mph",
}

response = requests.get(url, params=params)
print("Status code:", response.status_code)
print("Full URL:", response.url)

data = response.json()          # turn the JSON text into a Python dictionary
now = data["current"]
print("Time (UTC):", now["time"])
print("Temperature:", now["temperature_2m"], "°F")
print("Wind speed:", now["wind_speed_10m"], "mph")
\end{minted}

Here is what is happening:
\begin{itemize}
\item \verb|requests.get| sends a GET request. The \verb|params| dictionary becomes the \newterm{query string}\index{query string},
the part of the URL after the \verb|?|, with \verb|&| between each \verb|name=value| pair. The program prints the
full URL so you can see it. (The \verb|%2C| is how a comma is written in a URL.)
\item \verb|response.status_code| is the HTTP status code. 200 means it worked.
\item \verb|response.json()| turns the JSON in the response body into a Python dictionary. JSON objects become
dictionaries, and JSON lists become Python lists, so you reach into the data with square brackets, just as you
would with any dictionary.
\end{itemize}

When you run it, you should see something like this (your numbers will be different, because the weather
changes):
\begin{verbatim}
Status code: 200
Full URL: https://api.open-meteo.com/v1/forecast?latitude=39.74&longitude=-104.99&current=temperature_2m%2Cwind_speed_10m&temperature_unit=fahrenheit&wind_speed_unit=mph
Time (UTC): 2026-10-22T15:00
Temperature: 59.4 °F
Wind speed: 3.2 mph
\end{verbatim}

You can paste the full URL into a web browser, too. The browser sends the same GET request and shows you the
raw JSON that Python received. Every weather app on every phone is doing something very much like this, many
times a day.

\begin{Exercise}[title={Which Method?}, label=webapi_methods]
Which HTTP method would an app most likely use for each action?
\begin{enumerate}
\item Showing your profile page
\item Creating a new account
\item Changing only your profile picture
\item Deleting a post
\end{enumerate}
\end{Exercise}
\begin{Answer}[ref=webapi_methods]
\begin{enumerate}
\item GET, because it only reads data.
\item POST, because it creates something new.
\item PATCH, because it changes one field and leaves the rest alone. (Some APIs use PUT and send the whole profile.)
\item DELETE.
\end{enumerate}
\end{Answer}

\begin{Exercise}[title={Reading the JSON}, label=webapi_json]
Open-Meteo sends back JSON that includes this:
\begin{verbatim}
"current": {"time": "2026-10-22T15:00", "temperature_2m": 59.4}
\end{verbatim}
After \verb|data = response.json()|, what Python expression gives the temperature? What type of Python value is
it?
\end{Exercise}
\begin{Answer}[ref=webapi_json]
\verb|data["current"]["temperature_2m"]|. It is a \verb|float| (59.4), because JSON numbers with a decimal point
become Python floats.
\end{Answer}

\begin{Exercise}[title={Your Town}, label=webapi_town]
Look up the latitude and longitude of your town (a map app will show them), and change \filename{weather.py} to
print your own weather.
\end{Exercise}
\begin{Answer}[ref=webapi_town]
Change the two numbers in the \verb|params| dictionary. For example, Chicago is at about latitude 41.88 and
longitude $-87.63$. Longitudes west of Greenwich, England, are negative, and latitudes south of the equator are
negative.
\end{Answer}

\begin{Exercise}[title={A Three-Day Forecast}, label=webapi_forecast]
Open-Meteo can also send a daily forecast. Change the parameters to
\begin{minted}{python}
params = {
    "latitude": 39.74,
    "longitude": -104.99,
    "daily": "temperature_2m_max,temperature_2m_min",
    "temperature_unit": "fahrenheit",
    "timezone": "America/Denver",
    "forecast_days": 3,
}
\end{minted}
and print the date, high, and low for each of the three days. (Print \verb|data["daily"]| first to see how the
data is organized.)
\end{Exercise}
\begin{Answer}[ref=webapi_forecast]
\verb|data["daily"]| holds three lists of the same length: the dates, the highs, and the lows. Python's
\verb|zip| walks through them together:
\begin{minted}{python}
daily = data["daily"]
for day, high, low in zip(daily["time"], daily["temperature_2m_max"],
                          daily["temperature_2m_min"]):
    print(day, "high", high, "low", low)
\end{minted}
It prints something like \verb|2026-10-22 high 77.7 low 55.9|, one line per day.
\end{Answer}

\begin{Exercise}[title={A Bad Request}, label=webapi_bad]
Change the latitude to 139.74 and run \filename{weather.py} again. What status code do you get? Print
\verb|response.json()| to see why. Why is this a 400 and not a 404?
\end{Exercise}
\begin{Answer}[ref=webapi_bad]
The status code is 400, and the JSON says
\verb|{'error': True, 'reason': 'Latitude must be in range of -90 to 90°. Given: 139.74.'}|. (The program then
crashes, because there is no \verb|"current"| key in an error response. Real programs check
\verb|response.status_code| before using the data.) It is a 400 because the request itself was bad: no place on
Earth has a latitude of 139.74. A 404 would mean the address you asked for does not exist on the server, but the
\verb|/v1/forecast| endpoint is there.
\end{Answer}

In the next chapter, Building a Web API, you will switch sides and write a server of your own.
